\chapter{Capitolo 2 - Dalla cosa alla singolarità: base ontologica OCT}

\section{1. Problema}

Il Capitolo 1 ha fissato una distinzione necessaria: validità sintattica non coincide con validità ordinativa. Questa distinzione, tuttavia, resta incompleta se non viene chiarita la domanda ontologica di base: \textbf{che cosa è un ``oggetto'' quando la teoria deve descrivere sistemi reali dinamici, storici e ad alta dipendenza relazionale?}

Nel lessico classico, l'oggetto categoriale è un nodo astratto: non richiede struttura interna esplicita, e viene compreso attraverso la rete di morfismi che lo collega agli altri oggetti. Questa scelta ha una forza enorme, perché massimizza generalità e trasferibilità. Ma quando il quadro viene applicato a sistemi operativi concreti, emerge un limite ricorrente: il nodo astratto tende a nascondere differenze interne che, nel tempo, decidono persistenza o collasso.

In altri termini: due ``cose'' possono occupare la stessa posizione in una struttura formale e produrre esiti radicalmente diversi in presenza di perturbazioni, vincoli di risorsa, attriti semantici o mismatch di scala. Il formalismo classico registra correttamente la forma della composizione; non è progettato, da solo, per decidere la fecondità reale della composizione.

La conseguenza scientifica è netta. Se vogliamo che OCT entri nella grammatica della scienza umana, non basta aggiungere metriche a valle; serve una riformulazione a monte del referente ontologico. Non partiamo più dalla ``cosa'' come unità estensionale implicita, ma dalla \textbf{singolarità} come unità relazionale-funzionale esplicitamente situata in un contesto osservativo \texttt{Omega}.

Con ``cosa'' intendiamo qui un'entità trattata prevalentemente come supporto di proprietà e appartenenze. Con ``singolarità'' intendiamo un nodo la cui identità dipende da:

\begin{enumerate}
\def\labelenumi{\arabic{enumi}.}
\tightlist
\item
  profilo relazionale attivo;
\item
  capacità di mantenere coerenza interna sotto trasformazione;
\item
  funzione emergente non riducibile alla somma dei componenti;
\item
  traiettoria storica nel tempo osservato.
\end{enumerate}

Il passaggio non è lessicale, ma metodologico. Nel regime estensionale, la domanda è ``che cos'è questa entità?''. Nel regime ordinativo, la domanda diventa ``come questa singolarità mantiene identità e funzione in relazione ad altre singolarità, in \texttt{Omega}, lungo una traiettoria?''.

\subsection{1.1 Perché la nozione di ``cosa'' diventa insufficiente}

La nozione di cosa non è falsa; è insufficiente in domini dove:

\begin{enumerate}
\def\labelenumi{\arabic{enumi}.}
\tightlist
\item
  le relazioni cambiano la natura operativa dei nodi;
\item
  la storia delle interazioni modifica la risposta futura;
\item
  la funzionalità dipende da configurazioni collettive, non da attributi isolati.
\end{enumerate}

In tali domini, due entità con uguale descrizione estensionale possono divergere in modo irreversibile. Esempio astratto: due sistemi con lo stesso inventario di componenti e la stessa topologia iniziale possono differire nella qualità delle interfacce, nella sequenza di attivazione o nel grado di allineamento semantico; il risultato è che uno converge e l'altro degenera.

Questa divergenza non viene catturata da un'ontologia che privilegia l'enumerazione di proprietà statiche. Serve un'ontologia che assuma la relazione come costitutiva e non accessoria.

\subsection{1.2 Vincolo epistemico: evitare metafore ontologiche vaghe}

Il rischio opposto è usare ``singolarità'' in modo poetico o metafisico. Questo capitolo impone il vincolo contrario: la singolarità è definita solo se accompagnata da operatori osservabili, criteri di identità e condizioni di composizione.

In OCT non è accettabile dire che una singolarità ``esiste'' senza dichiarare:

\begin{enumerate}
\def\labelenumi{\arabic{enumi}.}
\tightlist
\item
  in quale \texttt{Omega} è valutata;
\item
  con quale metrica di coerenza;
\item
  con quale misura di funzione emergente;
\item
  con quale regime di stabilità temporale.
\end{enumerate}

Questa disciplina impedisce derive retoriche e tiene la teoria in un perimetro scientificamente falsificabile.

\begin{center}\rule{0.5\linewidth}{0.5pt}\end{center}

\section{2. Tesi del capitolo}

La tesi del capitolo è articolata in sei enunciati coordinati.

\begin{enumerate}
\def\labelenumi{\arabic{enumi}.}
\tightlist
\item
  \textbf{Tesi di rifondazione del referente}: l'unità ontologica primaria di OCT non è la cosa estensionale, ma la singolarità relazionale-funzionale.
\item
  \textbf{Tesi di identità situata}: l'identità di una singolarità non è assoluta; è definita rispetto a un contesto osservativo \texttt{Omega} e a una finestra temporale dichiarata.
\item
  \textbf{Tesi di irriducibilità funzionale}: una singolarità è scientificamente rilevante solo se mostra \texttt{Phi\_Omega\ !=\ 0}, cioè emergenza non banale.
\item
  \textbf{Tesi di persistenza}: la continuità identitaria richiede coerenza minima non degenerativa (\texttt{Coh\_Omega\ \textgreater{}=\ tau}) sotto trasformazioni ammissibili.
\item
  \textbf{Tesi di composizione responsabile}: composizioni sintatticamente corrette ma ordinativamente degradanti non devono essere trattate come equivalenti alle composizioni feconde.
\item
  \textbf{Tesi di continuità con il classico}: il passaggio a singolarità non abroga l'oggetto categoriale classico; ne introduce una interpretazione tipizzata più ricca per domini reali complessi.
\end{enumerate}

Queste tesi preparano il terreno per il Capitolo 3, in cui la categoria ordinativa verrà definita assiomaticamente (O1-O7). Senza una base ontologica esplicita, gli assiomi resterebbero sintatticamente eleganti ma semanticamente indeterminati.

\begin{center}\rule{0.5\linewidth}{0.5pt}\end{center}

\section{3. Definizioni canoniche}

Per rendere operativo il passaggio ontologico, introduciamo un nucleo definitorio minimo.

\subsection{Definizione 1 (Singolarità ordinativa)}

Una \textbf{singolarità ordinativa} \texttt{sigma} in contesto \texttt{Omega} è una tupla tipizzata:

\texttt{sigma\ :=\ \textless{}x,\ R\_Omega(x),\ H\_t(x),\ Pi\_Omega(x)\textgreater{}}

dove:

\begin{itemize}
\tightlist
\item
  \texttt{x} è il supporto identificativo minimo;
\item
  \texttt{R\_Omega(x)} è il profilo relazionale attivo in \texttt{Omega};
\item
  \texttt{H\_t(x)} è la traccia storica osservata in intervallo \texttt{t};
\item
  \texttt{Pi\_Omega(x)} è il potenziale ordinativo corrente.
\end{itemize}

Commento operativo:

\begin{itemize}
\tightlist
\item
  la singolarità non si riduce a \texttt{x};
\item
  \texttt{x} senza \texttt{R}, \texttt{H}, \texttt{Pi} è solo un indice tecnico.
\end{itemize}

\subsection{Definizione 2 (Profilo relazionale attivo)}

\texttt{R\_Omega(x)} è l'insieme tipizzato delle relazioni effettivamente operative per \texttt{x} in \texttt{Omega}, con pesi/qualità associati.

Commento:

\begin{itemize}
\tightlist
\item
  non tutte le relazioni possibili contano; solo quelle attive nel protocollo osservativo;
\item
  il profilo può cambiare con scala, regime temporale e perturbazioni.
\end{itemize}

\subsection{Definizione 3 (Potenziale ordinativo)}

\texttt{Pi\_Omega(x)} è la capacità di \texttt{x} di partecipare a composizioni che mantengono \texttt{Coh\_Omega} sopra soglia e generano \texttt{Phi\_Omega\ !=\ 0}.

Commento:

\begin{itemize}
\tightlist
\item
  \texttt{Pi} non misura prestazione assoluta;
\item
  misura predisposizione alla fecondità compositiva in un dominio dato.
\end{itemize}

\subsection{Definizione 4 (Identità situata)}

Due istanze \texttt{x\_a}, \texttt{x\_b} condividono la stessa identità situata in \texttt{Omega}, denotata \texttt{x\_a\ \textasciitilde{}\_Omega\ x\_b}, se e solo se:

\begin{enumerate}
\def\labelenumi{\arabic{enumi}.}
\tightlist
\item
  appartengono alla stessa classe tipologica minima;
\item
  preservano invarianti relazionali critici;
\item
  mantengono funzione emergente equivalente entro tolleranza dichiarata;
\item
  non mostrano divergenza degenerativa oltre soglia in finestra \texttt{t}.
\end{enumerate}

Commento:

\begin{itemize}
\tightlist
\item
  \texttt{\textasciitilde{}\_Omega} è più forte della mera equivalenza di etichetta;
\item
  è più debole della identità assoluta metafisica.
\end{itemize}

\subsection{Definizione 5 (Soglia di persistenza)}

Dato \texttt{tau\ in\ (0,1{]}}, una singolarità \texttt{sigma} è persistente in \texttt{Omega} su intervallo \texttt{t} se:

\texttt{inf\_t\ Coh\_Omega(sigma,\ t)\ \textgreater{}=\ tau}

e

\texttt{Phi\_Omega(sigma,\ t)} non collassa a zero su tutto l'intervallo.

Commento:

\begin{itemize}
\tightlist
\item
  la persistenza non richiede immobilità;
\item
  richiede stabilità non degenerativa.
\end{itemize}

\subsection{Definizione 6 (Degenerazione identitaria)}

Si ha degenerazione identitaria quando una singolarità conserva etichetta esterna ma perde le condizioni minime di coerenza e funzione che ne sostenevano il ruolo in \texttt{Omega}.

Indicatore minimo:

\texttt{Deg\_id\_Omega(sigma)\ :=\ 1} se \texttt{(Coh\ \textless{}\ tau)\ OR\ (Phi\ =\ 0)} su finestra critica.

\subsection{Definizione 7 (Morfismo ordinativamente ammissibile tra singolarità)}

Un morfismo \texttt{f:\ sigma\_i\ -\textgreater{}\ sigma\_j} è ordinativamente ammissibile se:

\begin{enumerate}
\def\labelenumi{\arabic{enumi}.}
\tightlist
\item
  è sintatticamente ben tipizzato;
\item
  non induce perdita di coerenza oltre soglia dichiarata;
\item
  preserva o incrementa \texttt{Phi\_Omega} in regime previsto;
\item
  rende esplicite le condizioni di fallimento.
\end{enumerate}

\begin{center}\rule{0.5\linewidth}{0.5pt}\end{center}

\section{4. Sviluppo formale}

\subsection{4.1 Dal nodo astratto al nodo tipizzato}

Nel lessico classico, un oggetto è spesso trattato come punto senza interno. Questo permette una notevole economia teorica. In OCT, per domini reali, adottiamo una rappresentazione intermedia:

\texttt{Obj\_OCT\ :=\ \textless{}carrier,\ rel\_profile,\ history,\ ord\_potential\textgreater{}}

La rappresentazione è minima ma sufficiente a evitare ambiguità ricorrenti:

\begin{enumerate}
\def\labelenumi{\arabic{enumi}.}
\tightlist
\item
  confondere identità nominale con persistenza funzionale;
\item
  confondere connettività con integrazione;
\item
  confondere esistenza formale con capacità generativa.
\end{enumerate}

Questa mossa non appesantisce inutilmente il formalismo: introduce solo i gradi di liberta necessari per rendere testabile la nozione di ``reale compositivo''.

\subsection{4.2 Criterio di realtà compositiva}

Sia \texttt{D} un diagramma di singolarità in \texttt{Omega}. Definiamo:

\texttt{Real\_OCT(D,\ Omega)\ :=\ Syntactic(D)\ AND\ Coh\_Omega(D)\ \textgreater{}=\ tau\ AND\ Phi\_Omega(D)\ !=\ 0}

Osservazione chiave:

\begin{itemize}
\tightlist
\item
  il criterio non elimina nessuna costruzione classica;
\item
  classifica quali costruzioni sono ontologicamente operative nel contesto.
\end{itemize}

Il criterio può essere raffinato con vincoli di robustezza:

\texttt{Robust\_OCT(D,\ Omega,\ Eps)\ :=\ Real\_OCT(D,\ Omega)\ AND\ Coh\_Omega(D\ +\ eps)\ \textgreater{}=\ tau\textquotesingle{}}

dove \texttt{eps} rappresenta perturbazioni compatibili con il dominio.

\subsection{4.3 Identità come traiettoria, non come istante}

Una fragilità dell'approccio estensionale applicato ai sistemi dinamici è la fissazione istantanea dell'identità. OCT assume invece un criterio di traiettoria:

\texttt{Id\_Omega(sigma)\ =\ stable\ if\ forall\ t\ in\ I,\ Coh\_Omega(sigma,t)\ \textgreater{}=\ tau\ and\ Phi\_Omega(sigma,t)\ sustains}

La forma precisa di ``sustains'' dipende dal dominio (continuità, media mobile, quantile minimo, ecc.), ma il principio resta:

\begin{enumerate}
\def\labelenumi{\arabic{enumi}.}
\tightlist
\item
  identità senza durata è un'etichetta;
\item
  identità con durata è una proprietà scientifica testabile.
\end{enumerate}

\subsection{4.4 Equivalenza classica e equivalenza ordinativa}

Siano \texttt{A}, \texttt{B} oggetti classici isomorfi (\texttt{A\ \textasciitilde{}=\ B} in senso strutturale). In OCT introduciamo:

\texttt{A\ \textasciitilde{}\_ord,Omega\ B} se le rispettive singolarizzazioni mantengono coerenza e funzione comparabili in \texttt{Omega}.

Ne segue:

\begin{enumerate}
\def\labelenumi{\arabic{enumi}.}
\tightlist
\item
  \texttt{A\ \textasciitilde{}=\ B} non implica necessariamente \texttt{A\ \textasciitilde{}\_ord,Omega\ B};
\item
  \texttt{A\ \textasciitilde{}\_ord,Omega\ B} richiede verifica empirico-formale contestualizzata.
\end{enumerate}

Questa asimmetria è il cuore dell'estensione ordinativa: preservare la correttezza algebrica e aggiungere discriminazione ontologica.

\subsection{4.5 Operatore di singolarizzazione}

Per connettere direttamente il classico con OCT introduciamo un operatore concettuale:

\texttt{S\_Omega:\ Obj\_class\ -\textgreater{}\ Sing\_OCT}

che associa a ogni oggetto classico \texttt{x} la singolarità \texttt{sigma} valutabile in \texttt{Omega}.

Requisiti minimi di \texttt{S\_Omega}:

\begin{enumerate}
\def\labelenumi{\arabic{enumi}.}
\tightlist
\item
  compatibilità tipologica (nessuna violazione di tipo);
\item
  esplicitazione del profilo relazionale rilevante;
\item
  dichiarazione della finestra storica osservata;
\item
  stima iniziale del potenziale ordinativo.
\end{enumerate}

\texttt{S\_Omega} non è unico in assoluto: varia con il protocollo osservativo. Questa dipendenza non è difetto, ma trasparenza metodologica.

\subsection{4.6 Regimi di perdita informativa}

Nel passaggio da sistemi reali a modellazione formale avvengono perdite di informazione. OCT distingue tre regimi:

\begin{enumerate}
\def\labelenumi{\arabic{enumi}.}
\tightlist
\item
  \textbf{perdita neutra}: riduzione che non altera coerenza/funzione rilevanti;
\item
  \textbf{perdita tollerabile}: riduzione che altera margini ma non cambia esito di validità;
\item
  \textbf{perdita critica}: riduzione che induce falsa equivalenza o falsa stabilità.
\end{enumerate}

Solo il terzo regime è inaccettabile scientificamente. Per questo ogni capitolo successivo dovrà dichiarare i passaggi di astrazione e il relativo rischio di perdita critica.

\subsection{4.7 Principio di non-indifferenza relazionale}

In ontologia ordinativa, non tutte le relazioni sono intercambiabili. Due architetture con stessa cardinalità relazionale possono avere qualità compositiva opposta.

Definiamo un principio metodologico:

\textbf{Principio NI (Non-Indifferenza)}\\
Se una variazione relazionale modifica sistematicamente \texttt{Coh\_Omega} o \texttt{Phi\_Omega}, tale variazione è ontologicamente rilevante e non può essere compressa in mera differenza notazionale.

Il principio NI protegge la teoria da semplificazioni che rendono invisibili i punti di rottura reali.

\subsection{4.8 Implicazione per la modellazione scientifica}

Con il passaggio alla singolarità:

\begin{enumerate}
\def\labelenumi{\arabic{enumi}.}
\tightlist
\item
  i dataset non sono solo collezioni di oggetti, ma campi di traiettorie relazionali;
\item
  i benchmark non valutano solo correttezza, ma tenuta ordinativa;
\item
  i protocolli devono includere stress test di coerenza e misura di emergenza.
\end{enumerate}

Questo spostamento è cruciale per la pubblicabilità interdisciplinare: rende esplicito come OCT traduce una posizione ontologica in una metodologia verificabile.

\begin{center}\rule{0.5\linewidth}{0.5pt}\end{center}

\section{5. Proposizioni e teoremi locali}

\subsection{Proposizione 2.1 (Insufficienza della sola identità estensionale)}

Enunciato: esistono coppie di sistemi \texttt{X}, \texttt{Y} con descrizione estensionale equivalente tali che \texttt{X} è ordinativamente persistente in \texttt{Omega} mentre \texttt{Y} non lo è.

Intuizione:

\begin{itemize}
\tightlist
\item
  equivalenza estensionale può ignorare differenze nel profilo relazionale attivo e nella traiettoria storica;
\item
  tali differenze sono sufficienti a produrre esiti divergenti su \texttt{Coh} e \texttt{Phi}.
\end{itemize}

Conseguenza:

\begin{itemize}
\tightlist
\item
  la sola identità estensionale non è criterio sufficiente di realtà compositiva.
\end{itemize}

Stato: \texttt{in\_proof} (programma di dimostrazione tramite controesempi costruttivi multi-dominio).

\subsection{Proposizione 2.2 (Necessità del profilo relazionale tipizzato)}

Enunciato: senza esplicitazione di \texttt{R\_Omega(x)}, la classificazione di validità ordinativa non è decidibile in generale.

Intuizione:

\begin{itemize}
\tightlist
\item
  \texttt{Coh\_Omega} dipende dalla qualità delle relazioni, non solo dalla presenza di connessioni;
\item
  senza \texttt{R\_Omega}, il valore di \texttt{Coh\_Omega} è sottodeterminato.
\end{itemize}

Conseguenza:

\begin{itemize}
\tightlist
\item
  ogni oggetto OCT deve essere accompagnato da un profilo relazionale minimo dichiarato.
\end{itemize}

Stato: \texttt{validated} a livello metodologico (\texttt{S1}), da elevare con formalizzazione completa nel Capitolo 9.

\subsection{Teorema 2.3 (Condizione minima di persistenza identitaria)}

Ipotesi:

\begin{enumerate}
\def\labelenumi{\arabic{enumi}.}
\tightlist
\item
  \texttt{sigma} è una singolarità in \texttt{Omega};
\item
  esiste \texttt{tau\ \textgreater{}\ 0} tale che \texttt{Coh\_Omega(sigma,t)\ \textgreater{}=\ tau} per ogni \texttt{t} in intervallo \texttt{I};
\item
  \texttt{Phi\_Omega(sigma,t)} non è identicamente nulla su \texttt{I}.
\end{enumerate}

Tesi:

\begin{itemize}
\tightlist
\item
  \texttt{sigma} mantiene identità ordinativa su \texttt{I} (nel senso di Definizione 5).
\end{itemize}

Proof sketch:

\begin{enumerate}
\def\labelenumi{\arabic{enumi}.}
\tightlist
\item
  dall'ipotesi 2 segue assenza di collasso coerentivo persistente;
\item
  dall'ipotesi 3 segue presenza di funzione emergente non triviale;
\item
  combinando 1 e 2 otteniamo le condizioni di persistenza definite in modo canonico.
\end{enumerate}

Conseguenze:

\begin{itemize}
\tightlist
\item
  la persistenza è verificabile con criteri minimi computabili;
\item
  l'identità ordinativa diventa proprietà operativa, non assunzione.
\end{itemize}

Stato: \texttt{in\_proof} (chiusura formale prevista in blocco teorematico F01-F10).

\subsection{Teorema 2.4 (Non-equivalenza tra isomorfismo e equivalenza ordinativa piena)}

Ipotesi:

\begin{enumerate}
\def\labelenumi{\arabic{enumi}.}
\tightlist
\item
  \texttt{A\ \textasciitilde{}=\ B} (isomorfismo classico);
\item
  esiste contesto \texttt{Omega} in cui le singolarizzazioni \texttt{S\_Omega(A)}, \texttt{S\_Omega(B)} mostrano profili relazionali non allineabili rispetto a criterio NI;
\item
  differenza di profilo induce \texttt{Coh} o \texttt{Phi} non equivalenti oltre soglia.
\end{enumerate}

Tesi:

\begin{itemize}
\tightlist
\item
  non vale \texttt{A\ \textasciitilde{}\_ord,Omega\ B}.
\end{itemize}

Proof sketch:

\begin{enumerate}
\def\labelenumi{\arabic{enumi}.}
\tightlist
\item
  l'isomorfismo garantisce intercambiabilità strutturale nel quadro classico;
\item
  il criterio ordinativo richiede anche allineamento su coerenza/funzione;
\item
  per ipotesi tale allineamento manca, quindi equivalenza ordinativa decade.
\end{enumerate}

Conseguenze:

\begin{itemize}
\tightlist
\item
  OCT è un'estensione conservativa ma non banale della teoria classica;
\item
  serve doppia verifica: algebrica e ordinativa.
\end{itemize}

Stato: \texttt{revise} (richiede formalizzazione completa di \texttt{\textasciitilde{}\_ord,Omega} nel Capitolo 3 e validazione nei casi studio).

\begin{center}\rule{0.5\linewidth}{0.5pt}\end{center}

\section{6. Implicazioni operative}

Le definizioni del capitolo hanno impatti immediati su ricerca, ingegneria e validazione.

\subsection{6.1 Progettazione di modelli}

Un modello OCT-compatibile deve dichiarare, per ogni unità analitica:

\begin{enumerate}
\def\labelenumi{\arabic{enumi}.}
\tightlist
\item
  supporto identificativo (\texttt{x});
\item
  profilo relazionale osservato (\texttt{R\_Omega});
\item
  finestra storica (\texttt{H\_t});
\item
  ipotesi sul potenziale ordinativo (\texttt{Pi\_Omega}).
\end{enumerate}

Senza questi quattro elementi, il modello resta nel regime descrittivo classico e non può sostenere claim ordinativi forti.

\subsection{6.2 Disegno dei benchmark}

I benchmark devono essere riformulati per catturare non solo accuratezza locale ma tenuta compositiva:

\begin{enumerate}
\def\labelenumi{\arabic{enumi}.}
\tightlist
\item
  test di coerenza sotto perturbazione;
\item
  misura della persistenza di funzione emergente;
\item
  tracciamento della degenerazione identitaria.
\end{enumerate}

In pratica, un benchmark OCT non chiede solo ``funziona?'', ma ``continua a funzionare senza collasso di coerenza quando il contesto cambia entro i limiti dichiarati?''.

\subsection{6.3 Pipeline di pubblicazione e audit}

Per pubblicabilità su GitHub/Zenodo/OSF e piattaforme analoghe, il capitolo suggerisce una disciplina minima:

\begin{enumerate}
\def\labelenumi{\arabic{enumi}.}
\tightlist
\item
  dichiarare il contesto \texttt{Omega} nei metadati di ogni esperimento;
\item
  allegare protocollo di misura per \texttt{Coh/Phi/Delta};
\item
  classificare i claim con \texttt{S0-S3};
\item
  includere limiti e failure mode espliciti.
\end{enumerate}

Questa disciplina riduce ambiguità interpretative e migliora replicabilità internazionale.

\subsection{6.4 Impatto su sistemi AI ordinativi}

Per agenti AI ispirati al framework TE/OCT, il passaggio da cosa a singolarità implica:

\begin{enumerate}
\def\labelenumi{\arabic{enumi}.}
\tightlist
\item
  rappresentazioni interne orientate a traiettorie e relazioni, non solo stati statici;
\item
  controllo decisionale che penalizza composizioni degradanti anche se formalmente lecite;
\item
  governance del ciclo di vita basata su segnali di degenerazione precoce.
\end{enumerate}

L'obiettivo non è antropomorfizzare il sistema, ma evitare modelli che ottimizzano metriche locali producendo collasso globale.

\subsection{6.5 Interoperabilità con la teoria classica}

Per evitare fratture inutili con la letteratura esistente, ogni costrutto OCT deve indicare:

\begin{enumerate}
\def\labelenumi{\arabic{enumi}.}
\tightlist
\item
  equivalente classico di partenza;
\item
  limite classico rilevante;
\item
  estensione ordinativa introdotta.
\end{enumerate}

Questa tripla è fondamentale per un dialogo serio con la comunità matematica e con le scienze applicate.

\begin{center}\rule{0.5\linewidth}{0.5pt}\end{center}

\section{7. Limiti e non-claim}

Questo capitolo introduce una base ontologica operativa, ma non pretende più di quanto può dimostrare in questa fase.

Limiti espliciti:

\begin{enumerate}
\def\labelenumi{\arabic{enumi}.}
\tightlist
\item
  non viene ancora fornita una assiomatizzazione completa della categoria ordinativa (rimandata al Capitolo 3);
\item
  non viene ancora chiusa una prova completa dei teoremi 2.3 e 2.4;
\item
  non è ancora fissata una metrica universale di \texttt{Phi\_Omega} valida per tutti i domini;
\item
  la nozione di identità situata dipende da soglie e protocolli dominio-specifici.
\end{enumerate}

Failure mode riconosciuti:

\begin{enumerate}
\def\labelenumi{\arabic{enumi}.}
\tightlist
\item
  soglie \texttt{tau} mal calibrate possono produrre falsi positivi/negativi di persistenza;
\item
  profili relazionali incompleti possono simulare stabilità inesistente;
\item
  finestre temporali troppo corte possono mascherare degenerazioni lente.
\end{enumerate}

Box C - Non-claim

Questo capitolo non implica:

\begin{enumerate}
\def\labelenumi{\arabic{enumi}.}
\tightlist
\item
  che la nozione classica di oggetto sia errata;
\item
  che ogni divergenza empirica richieda nuova ontologia;
\item
  che OCT abbia già chiuso tutte le prove formali necessarie.
\end{enumerate}

\begin{center}\rule{0.5\linewidth}{0.5pt}\end{center}

\section{8. Claim Ledger (S0-S3)}

{\def\LTcaptype{none} % do not increment counter
\begin{longtable}[]{@{}lllll@{}}
\toprule\noalign{}
ID & Claim & Tipo & Evidenza & Stato \\
\midrule\noalign{}
\endhead
\bottomrule\noalign{}
\endlastfoot
C2.1 & La nozione estensionale di ``cosa'' è insufficiente per discriminare persistenza in sistemi dinamici complessi & interpretativo-metodologico & S1 & validated \\
C2.2 & L'unità primaria OCT deve essere la singolarità \texttt{sigma\ :=\ \textless{}x,R,H,Pi\textgreater{}} & formale-programmatico & S2 & revise \\
C2.3 & Senza \texttt{R\_Omega} la validità ordinativa non è decidibile in generale & formale & S1 & validated \\
C2.4 & Isomorfismo classico non implica equivalenza ordinativa piena in \texttt{Omega} & formale & S2 & in\_review \\
C2.5 & Persistenza identitaria richiede congiuntamente soglia di coerenza e funzione emergente non nulla & formale & S2 & in\_proof \\
C2.6 & Benchmark OCT devono includere stress test di coerenza e tracciamento degenerativo & operativo & S1 & validated \\
C2.7 & Il passaggio alla singolarità migliora la replicabilità inter-dominio se accompagnato da metadati \texttt{Omega} espliciti & metodologico & S1 & in\_review \\
C2.8 & Non esiste, in questo stadio, una misura unica universale di \texttt{Phi\_Omega} & epistemico & S0 & validated \\
\end{longtable}
}

\begin{center}\rule{0.5\linewidth}{0.5pt}\end{center}

\section{9. Bridge al Capitolo 3}

Il Capitolo 2 ha stabilito il referente ontologico minimo di OCT: non la cosa estensionale, ma la singolarità relazionale-funzionale situata in \texttt{Omega}.

Con questa base possiamo compiere il passaggio successivo: formulare la \textbf{Categoria Ordinativa} in senso stretto. Il Capitolo 3 introdurrà:

\begin{enumerate}
\def\labelenumi{\arabic{enumi}.}
\tightlist
\item
  dominio degli oggetti-singolarità e dei morfismi ordinativamente ammissibili;
\item
  assiomi O1-O7 per composizione, identità, coerenza, emergenza e non-degenerazione;
\item
  prime conseguenze strutturali per limiti, equivalenze e funtorialità in quadro OCT.
\end{enumerate}

In termini operativi, il passaggio è questo:

\begin{itemize}
\tightlist
\item
  Capitolo 1: perché la validità classica non basta nei domini reali;
\item
  Capitolo 2: quale ontologia minima serve per colmare il gap;
\item
  Capitolo 3: come trasformare questa ontologia in teoria categoriale formalizzata.
\end{itemize}

\begin{center}\rule{0.5\linewidth}{0.5pt}\end{center}

\section{Riferimenti}

Riferimenti di framework interno:

\begin{enumerate}
\def\labelenumi{\arabic{enumi}.}
\tightlist
\item
  Ghioni, F. (2026). \texttt{TE\_CORE\_v5.1.md}.
\item
  Ghioni, F. (2026). \texttt{Ordinative\_Set\_Theory\_OST\_A\_Concise\_Guide\_For\_AI\_v2\_1.md}.
\item
  \texttt{OCT\_CLASSICAL\_TO\_ORDINATIVE\_MAP\_v0\_1.md}.
\item
  \texttt{OCT\_THEOREM\_PROGRAM\_v0\_1.md}.
\item
  \texttt{OCT\_TYPED\_FORMAL\_SPEC\_v0\_1.md}.
\end{enumerate}

Riferimenti primari (contesto disciplinare):

\begin{enumerate}
\def\labelenumi{\arabic{enumi}.}
\tightlist
\item
  Eilenberg, S., Mac Lane, S. (1945). General Theory of Natural Equivalences.
\item
  Mac Lane, S. (1998). Categories for the Working Mathematician.
\item
  Spivak, D. (2014). Category Theory for the Sciences.
\end{enumerate}
